\documentclass[10pt,a4paper]{amsart}
\usepackage{amsmath,amssymb,amsthm,graphicx}
\usepackage{verbatim}
\newenvironment{reference}{\comment}{\endcomment}
%\newenvironment{reference}{}{}
\newenvironment{slogan}{\comment}{\endcomment}
\newenvironment{history}{\comment}{\endcomment}

% For commutative diagrams we use Xy-pic
\usepackage[all]{xy}

% We use 2cell for 2-commutative diagrams.
\xyoption{2cell}
% Theorem environments.
%
\theoremstyle{plain}
\newtheorem{theorem}[subsection]{Theorem}
\newtheorem{proposition}[subsection]{Proposition}
\newtheorem{lemma}[subsection]{Lemma}

\theoremstyle{definition}
\newtheorem{definition}[subsection]{Definition}
\newtheorem{example}[subsection]{Example}
\newtheorem{exercise}[subsection]{Exercise}
\newtheorem{situation}[subsection]{Situation}

\theoremstyle{remark}
\newtheorem{remark}[subsection]{Remark}
\newtheorem{remarks}[subsection]{Remarks}

\numberwithin{equation}{subsection}

% Macros
%
\def\lim{\mathop{\mathrm{lim}}\nolimits}
\def\colim{\mathop{\mathrm{colim}}\nolimits}
\def\Spec{\mathop{\mathrm{Spec}}}
\def\Hom{\mathop{\mathrm{Hom}}\nolimits}
\def\Ext{\mathop{\mathrm{Ext}}\nolimits}
\def\SheafHom{\mathop{\mathcal{H}\!\mathit{om}}\nolimits}
\def\SheafExt{\mathop{\mathcal{E}\!\mathit{xt}}\nolimits}
\def\Sch{\mathit{Sch}}
\def\Mor{\mathop{\mathrm{Mor}}\nolimits}
\def\Ob{\mathop{\mathrm{Ob}}\nolimits}
\def\Sh{\mathop{\mathit{Sh}}\nolimits}
\def\NL{\mathop{N\!L}\nolimits}
\def\CH{\mathop{\mathrm{CH}}\nolimits}
\def\proetale{{pro\text{-}\acute{e}tale}}
\def\etale{{\acute{e}tale}}
\def\QCoh{\mathit{QCoh}}
\def\Ker{\mathop{\mathrm{Ker}}}
\def\Im{\mathop{\mathrm{Im}}}
\def\Coker{\mathop{\mathrm{Coker}}}
\def\Coim{\mathop{\mathrm{Coim}}}

% Boxtimes
%
\DeclareMathSymbol{\boxtimes}{\mathbin}{AMSa}{"02}

%
% Macros for moduli stacks/spaces
%
\def\QCohstack{\mathcal{QC}\!\mathit{oh}}
\def\Cohstack{\mathcal{C}\!\mathit{oh}}
\def\Spacesstack{\mathcal{S}\!\mathit{paces}}
\def\Quotfunctor{\mathrm{Quot}}
\def\Hilbfunctor{\mathrm{Hilb}}
\def\Curvesstack{\mathcal{C}\!\mathit{urves}}
\def\Polarizedstack{\mathcal{P}\!\mathit{olarized}}
\def\Complexesstack{\mathcal{C}\!\mathit{omplexes}}
% \Pic is the operator that assigns to X its picard group, usage \Pic(X)
% \Picardstack_{X/B} denotes the Picard stack of X over B
% \Picardfunctor_{X/B} denotes the Picard functor of X over B
\def\Pic{\mathop{\mathrm{Pic}}\nolimits}
\def\Picardstack{\mathcal{P}\!\mathit{ic}}
\def\Picardfunctor{\mathrm{Pic}}
\def\Deformationcategory{\mathcal{D}\!\mathit{ef}}
\usepackage{cleveref}
\crefname{originallemma}{original lemma}{original lemmas}
\Crefname{originallemma}{Original lemma}{Original lemmas}
\makeatletter
\newlabel{divisors-lemma-blowing-up-projective}{{02NS}{source}{}{sourcepointer-1}{}}
\newlabel{divisors-lemma-blowing-up-projective@cref}{{[originallemma][0][]02NS}{[1][0][]source}{}{}{}}
\makeatother
\crefname{originallemma}{original lemma}{original lemmas}
\Crefname{originallemma}{Original lemma}{Original lemmas}
\makeatletter
\newlabel{divisors-lemma-pullback-effective-Cartier-defined}{{02OO}{source}{}{sourcepointer-2}{}}
\newlabel{divisors-lemma-pullback-effective-Cartier-defined@cref}{{[originallemma][0][]02OO}{[1][0][]source}{}{}{}}
\makeatother
\crefname{originalreference}{original reference}{original references}
\Crefname{originalreference}{Original reference}{Original references}
\makeatletter
\newlabel{lemma-sections-powers-I-rational}{{0B4Y}{source}{}{sourcepointer-3}{}}
\newlabel{lemma-sections-powers-I-rational@cref}{{[originalreference][0][]0B4Y}{[1][0][]source}{}{}{}}
\makeatother
\crefname{originallemma}{original lemma}{original lemmas}
\Crefname{originallemma}{Original lemma}{Original lemmas}
\makeatletter
\newlabel{morphisms-lemma-nagata-normalization}{{035S}{source}{}{sourcepointer-4}{}}
\newlabel{morphisms-lemma-nagata-normalization@cref}{{[originallemma][0][]035S}{[1][0][]source}{}{}{}}
\makeatother
\crefname{originallemma}{original lemma}{original lemmas}
\Crefname{originallemma}{Original lemma}{Original lemmas}
\makeatletter
\newlabel{properties-lemma-ample-quasi-coherent}{{01QJ}{source}{}{sourcepointer-5}{}}
\newlabel{properties-lemma-ample-quasi-coherent@cref}{{[originallemma][0][]01QJ}{[1][0][]source}{}{}{}}
\makeatother
\crefname{originalreference}{original reference}{original references}
\Crefname{originalreference}{Original reference}{Original references}
\makeatletter
\newlabel{situation-rational}{{0B4W}{source}{}{sourcepointer-6}{}}
\newlabel{situation-rational@cref}{{[originalreference][0][]0B4W}{[1][0][]source}{}{}{}}
\makeatother
\title{The blowup of the maximal ideal of a normal surface singularity defining a rational singularity is normal}
\author{The Stacks Project Authors}
\date{stacks-project:resolve.tex:local-snapshot-20261012; GNU FDL 1.2 or later}
\begin{document}
\maketitle
\noindent\textit{Editorial note.} All delivered mathematical text is the author's own, copied from the pinned Stacks Project chapter file resolve.tex; no statement or proof step was rewritten, reordered or removed.\par\medskip
\noindent\textit{Editorial note.} The counted proof is the author's complete argument: the description of the blowup through the Rees algebra with the pullback of the exceptional divisor along the finite normalization, the reduction to an equality of sections of all twists through the ample invertible sheaf, and the computation of both sides using the rational singularity input and torsion freeness.\par\medskip
\noindent\textit{Editorial note.} The label lemma-blow-up-normal-rational names the counted statement, and the proof environment immediately following it in the source is delivered verbatim as its proof.\par\medskip
\noindent\textit{Editorial note.} The counted statement is carried by the author's own lemma environment in the chapter source.\par\medskip
\noindent\textit{Editorial note.} Pre-open gate: python operations/agent-authoring/tools/precard\_gate.py --source resolve.tex --label lemma-blow-up-normal-rational --start 2533 --end 2578 --local-id 28093 returned RESULT=OK (L28093 free, occupancy FREE, public interval clear). Pre-assembly tag self-check: canonical tag resolved and not present in reference/public-items.jsonl.\par\medskip
\noindent\textit{Editorial note.} The statement and proof are copied from the pinned chapter file; the cited inputs are kept as external pointers with their canonical Stacks tags.\par\medskip
\noindent\textit{Allowed standard prerequisites (editorial).}\begin{itemize}
\item Description of the blowup as the projective spectrum of the Rees algebra. Divisors, the lemma identifying the blowup with the projective spectrum of the Rees algebra and the lemma identifying the ideal sheaf of the exceptional divisor with the twisting sheaf. Supplies the integrality of the blowup, its projectivity over the base and the identification of the ideal sheaf of the exceptional divisor.
\item Finiteness of normalization for Nagata rings and pullback of effective Cartier divisors. Morphisms, the lemma that normalization of a Nagata scheme is finite, and Divisors, the lemma that the pullback of an effective Cartier divisor along a finite morphism is defined. Supply the finiteness of the normalization and the fact that the inverse image of the exceptional divisor is an effective Cartier divisor.
\item Recovery of a module from its associated graded along an ample invertible sheaf. Properties, the lemma that a quasi-coherent module can be recovered from its associated graded module with respect to an ample invertible sheaf. Supplies the criterion that if the sections of all twists agree then the modules agree.
\item Sections of the powers of the ideal on the normalized blowup. Same chapter, the lemma computing the sections of the powers of the ideal sheaf on the normalized blowup in the rational singularity situation. Supplies the identification of the sections of the pullback of the twisting sheaf with the powers of the maximal ideal.
\end{itemize}\medskip
\ifdefined\hypertarget\hypertarget{sourcepointer-1}{}\fi
\noindent\textit{Source reference.} Original lemma 02NS: sources/stacks-chow/tags.txt, source lines 3491--3491. Prerequisite pointer invoked at proof line(s) 2553; it names a standard fact of a sibling chapter that the counted argument uses, and every step of the counted proof built on it is retained verbatim in the delivered body.\par\medskip
\ifdefined\hypertarget\hypertarget{sourcepointer-2}{}\fi
\noindent\textit{Source reference.} Original lemma 02OO: sources/stacks-chow/tags.txt, source lines 3523--3523. Prerequisite pointer invoked at proof line(s) 2556; it names a standard fact of a sibling chapter that the counted argument uses, and every step of the counted proof built on it is retained verbatim in the delivered body.\par\medskip
\ifdefined\hypertarget\hypertarget{sourcepointer-3}{}\fi
\noindent\textit{Source reference.} Original reference 0B4Y: sources/stacks-resolve/pinned-original-source.tex, source lines 2488--2488. Prerequisite pointer invoked at proof line(s) 2573; it names a standard fact of the same chapter that the counted argument uses, and every step of the counted proof built on it is retained verbatim in the delivered body.\par\medskip
\ifdefined\hypertarget\hypertarget{sourcepointer-4}{}\fi
\noindent\textit{Source reference.} Original lemma 035S: sources/stacks-chow/tags.txt, source lines 4139--4139. Prerequisite pointer invoked at proof line(s) 2548; it names a standard fact of a sibling chapter that the counted argument uses, and every step of the counted proof built on it is retained verbatim in the delivered body.\par\medskip
\ifdefined\hypertarget\hypertarget{sourcepointer-5}{}\fi
\noindent\textit{Source reference.} Original lemma 01QJ: sources/stacks-chow/tags.txt, source lines 2268--2268. Prerequisite pointer invoked at proof line(s) 2566; it names a standard fact of a sibling chapter that the counted argument uses, and every step of the counted proof built on it is retained verbatim in the delivered body.\par\medskip
\ifdefined\hypertarget\hypertarget{sourcepointer-6}{}\fi
\noindent\textit{Source reference.} Original reference 0B4W: sources/stacks-resolve/pinned-original-source.tex, source lines 2455--2455. Prerequisite pointer invoked at proof line(s) 2535; it names a standard fact of the same chapter that the counted argument uses, and every step of the counted proof built on it is retained verbatim in the delivered body.\par\medskip
\begin{lemma}
\label{lemma-blow-up-normal-rational}
In Situation \ref{situation-rational}
the blowup of $\Spec(A)$ in $\mathfrak m$ is normal.
\end{lemma}

\begin{proof}
Let $X' \to \Spec(A)$ be the blowup, in other words
$$
X' = \text{Proj}(A \oplus \mathfrak m \oplus \mathfrak m^2 \oplus \ldots).
$$
is the Proj of the Rees algebra. This in particular shows that
$X'$ is integral and that $X' \to \Spec(A)$ is a projective
modification. Let $X$ be the normalization of $X'$.
Since $A$ is Nagata, we see that $\nu : X \to X'$ is finite
(Morphisms, Lemma \ref{morphisms-lemma-nagata-normalization}).
Let $E' \subset X'$ be the exceptional divisor and let $E \subset X$
be the inverse image. Let $\mathcal{I}' \subset \mathcal{O}_{X'}$
and $\mathcal{I} \subset \mathcal{O}_X$ be their ideal sheaves.
Recall that $\mathcal{I}' = \mathcal{O}_{X'}(1)$
(Divisors, Lemma \ref{divisors-lemma-blowing-up-projective}).
Observe that $\mathcal{I} = \nu^*\mathcal{I}'$ and that $E$ is an
effective Cartier divisor (Divisors, Lemma
\ref{divisors-lemma-pullback-effective-Cartier-defined}).
We are trying to show that $\nu$ is an isomorphism. As $\nu$ is finite,
it suffices to show that $\mathcal{O}_{X'} \to \nu_*\mathcal{O}_X$
is an isomorphism. If not, then we can find an $n \geq 0$ such that
$$
H^0(X', (\mathcal{I}')^n) \not =
H^0(X', (\nu_*\mathcal{O}_X) \otimes (\mathcal{I}')^n)
$$
for example because we can recover quasi-coherent $\mathcal{O}_{X'}$-modules
from their associated graded modules, see
Properties, Lemma \ref{properties-lemma-ample-quasi-coherent}.
By the projection formula we have
$$
H^0(X', (\nu_*\mathcal{O}_X) \otimes (\mathcal{I}')^n) =
H^0(X, \nu^*(\mathcal{I}')^n) =
H^0(X, \mathcal{I}^n) = \mathfrak m^n
$$
the last equality by Lemma \ref{lemma-sections-powers-I-rational}.
On the other hand, there is clearly an injection
$\mathfrak m^n \to H^0(X', (\mathcal{I}')^n)$. Since
$H^0(X', (\mathcal{I}')^n)$ is torsion free we conclude equality holds
for all $n$, hence $X = X'$.
\end{proof}

\end{document}
